\documentclass[11pt]{article}
\usepackage[margin=1.1in]{geometry}
\usepackage{amsmath,amssymb,amsthm}
\usepackage{booktabs}
\usepackage{microtype}
\usepackage[colorlinks=true,linkcolor=blue,citecolor=blue]{hyperref}

\newtheorem{proposition}{Proposition}
\newtheorem{lemma}{Lemma}
\theoremstyle{remark}
\newtheorem{remark}{Remark}

\newcommand{\dd}{\,\mathrm{d}}
\newcommand{\code}[1]{\texttt{#1}}
\newcommand{\eps}{\varepsilon_{\mathrm{m}}}
\DeclareMathOperator{\sgn}{sgn}

\title{Bessel functions past the edge of the format\\
\large The evaluation paths in \texttt{chladni.helmholtz.\_bessel}}
\author{chladni development notes}
\date{}

\begin{document}
\maketitle

\begin{abstract}
\noindent
A corner-adapted Fourier--Bessel basis needs $J_\nu(kr)$ at orders $\nu$ in the hundreds
with $kr$ small, and in that regime \code{scipy.special.jv} returns exactly zero while the
true value is finite. This note develops what the code does instead: a factorisation
$J_\nu(x)=S\,e^{\sigma}$ with an \emph{analytic} scale $\sigma$, three evaluation branches
for $S$, and the argument that the branches overlap rather than abut. It also derives the
error floor --- which turns out to be $\sigma$ itself and not the summation --- and the
separate per-column scale that leaves the basis matrix well conditioned. Every numerical
claim here is measured --- the accuracy figures against \code{mpmath} at 50 digits --- and
reproducible from the code.
\end{abstract}

\section{The problem}

The method of particular solutions expands an eigenfunction near a corner of interior angle
$\omega$ in the functions $J_{\mu}(kr)\sin(\mu\theta)$ with $\mu=\ell\nu$ and
$\nu=\pi/\omega$. A sharp corner makes $\nu$ large, so a basis of a dozen degrees reaches
$\mu\sim 300$; meanwhile $r$ is bounded by the domain and $k$ by the eigenvalue sought, so
the argument $x=kr$ stays modest. The interesting regime is therefore $\nu\gg x$, where
$J_\nu(x)$ is very small.

How small is the difficulty. Table~\ref{tab:cliff} bisects the point at which
\code{scipy.special.jv} first returns a nonzero value. The behaviour is not a gradual loss
of accuracy: the routine holds thirteen to fourteen digits and then flushes to exactly
$0.0$, at $|J|\approx10^{-291}$ --- some seventeen decades above the smallest normal double.

\begin{table}[h]
\centering
\begin{tabular}{lrrrrrr}
\toprule
$\nu$ & $100$ & $200$ & $300$ & $400$ & $500$ & $700$\\
\midrule
smallest $x$ with $\code{jv}(\nu,x)\neq 0$ & $0.096$ & $5.32$ & $23.98$ & $51.74$ & $92.06$ & $194.1$\\
\bottomrule
\end{tabular}
\caption{The underflow edge, bisected against \code{scipy.special.jv} itself.}
\label{tab:cliff}
\end{table}

A column of zeros is not a small column; it is no column at all. The predecessor code,
\code{lappy}, evaluated exactly this way and its benchmark notes concluded that the sharp
corner columns ``carry almost no independent information about the domain''. Half of that is
true --- near-dependence at a sharp corner is intrinsic to double precision --- and half is
an artefact of the representation, which is what this note removes.

\section{The factorisation}

Write
\begin{equation}
  \sigma_J(\nu,x) \;=\; \nu\ln\frac{x}{2} \;-\; \ln\Gamma(\nu+1),
  \qquad
  J_\nu(x) \;=\; S\,e^{\sigma_J},
  \qquad
  J_\nu'(x) \;=\; S'\,e^{\sigma_J}.
  \label{eq:factor}
\end{equation}
The quantity $e^{\sigma_J}=(x/2)^\nu/\Gamma(\nu+1)$ is the leading term of the ascending
series. Two properties make this the right choice of scale rather than one of many.

\begin{proposition}[$|S|\le 1$]
For real $x\ge 0$ and $\nu\ge -\tfrac12$,
$\;|J_\nu(x)| \le (x/2)^\nu/\Gamma(\nu+1)$, so $|S|\le 1$ identically.
\end{proposition}

\begin{proof}
Poisson's integral representation gives
\[
  J_\nu(x)=\frac{(x/2)^\nu}{\sqrt{\pi}\,\Gamma(\nu+\tfrac12)}
           \int_{-1}^{1}(1-t^2)^{\nu-1/2}\cos(xt)\dd t ,
\]
valid for $\nu>-\tfrac12$. Bounding $|\cos(xt)|\le 1$ and using
$\int_{-1}^{1}(1-t^2)^{\nu-1/2}\dd t=\sqrt{\pi}\,\Gamma(\nu+\tfrac12)/\Gamma(\nu+1)$
gives the claim.
\end{proof}

So $|S|\le 1$ is a structural invariant, independent of which branch computed it, and the
test suite asserts it directly. It also means $\sigma_J$ is an \emph{upper} scale: the
factorisation never has to represent a number larger than one.

The second property is that $\sigma_J$ is analytic and depends on the evaluation point only
through $x$ --- it is never a logarithm of the computed value. That matters in the
oscillatory regime $\nu\lesssim x$, where $\ln|J_\nu|$ is genuinely singular at every zero.
A scale defined as $\ln|J|$ would inherit those poles; this one does not, and the scaled
value $S$ carries the oscillation, the sign changes and the zeros itself.

Finally, because $\sigma_J$ is the same expression in every branch, a caller may subtract
two of them --- which is exactly what the basis does when it converts from the pointwise
scale to a per-column one (\S\ref{sec:column}) --- without knowing how either was obtained.

\section{Branch A: the ascending series}

For $x$ not too large relative to $\sqrt{\nu}$ the defining series is both convergent and
accurate. From
\[
  J_\nu(x)=\sum_{j\ge0}\frac{(-1)^j (x/2)^{\nu+2j}}{j!\,\Gamma(\nu+j+1)}
\]
and~\eqref{eq:factor},
\begin{equation}
  S=\sum_{j\ge0} t_j,
  \qquad t_0=1,
  \qquad t_j=t_{j-1}\,\frac{q}{j(\nu+j)},
  \qquad q=-\frac{x^2}{4}.
  \label{eq:series}
\end{equation}

\subsection{The derivative comes free}

Each term of the series is $c_j (x/2)^{\nu+2j}$ with $c_j$ independent of $x$, so
\[
  \frac{\dd}{\dd x}\Big[(x/2)^{\nu+2j}\Big]
  = \frac{\nu+2j}{x}\,(x/2)^{\nu+2j},
\]
and therefore
\begin{equation}
  x\,S' \;=\; \sum_{j\ge0} (\nu+2j)\,t_j .
  \label{eq:deriv}
\end{equation}
The same $t_j$ serve both sums, so one loop with two accumulators yields the value and the
derivative. This is worth stating because the alternative --- forming $J_\nu'$ from
$\tfrac12(J_{\nu-1}-J_{\nu+1})$ --- would need two more series \emph{and} would produce them
with two different scales $\sigma_J(\nu\mp1,x)$, which then have to be reconciled.

\subsection{Accuracy, and the largest term}

Equation~\eqref{eq:series} is an alternating sum. Floating-point summation commits an
absolute error of order $\eps M$ where $M=\max_j|t_j|$, so the relative error of $S$ is about
$\eps M/|S|$. The useful question is how big $M$ gets.

For $j\ll\nu$ the recurrence ratio is
\[
  \left|\frac{t_j}{t_{j-1}}\right| = \frac{x^2/4}{j(\nu+j)} \approx \frac{1}{j}\cdot\frac{x^2}{4\nu},
\]
which is the ratio of the exponential series $\sum_j \lambda^j/j!$ with $\lambda=x^2/(4\nu)$.
Hence
\begin{equation}
  M \;\lesssim\; \exp\!\left(\frac{x^2}{4\nu}\right),
  \label{eq:maxterm}
\end{equation}
attained near $j\approx\lambda$. The estimate is conservative --- it ignores both the
$j$-dependence of $\nu+j$ and the $1/\sqrt{2\pi\lambda}$ in Stirling's approximation of the
largest term --- and measurement confirms it errs on the safe side by a factor of a few.
The code uses~\eqref{eq:maxterm} with $\nu+1$ in place of $\nu$, which keeps $\nu=0$ finite
and errs further toward the library branch, where the library is exact anyway.

Table~\ref{tab:seriesacc} shows the measured relative error against \code{mpmath} at
50 digits, alongside the predicted $M$. The error tracks $\eps M$ closely, and the threshold
the code uses is $M\le10^3$, equivalently $x\lesssim 5.25\sqrt{\nu}$.

\begin{table}[h]
\centering
\begin{tabular}{rrrrr}
\toprule
$\nu$ & $x$ & $x/\sqrt{\nu}$ & $M\approx e^{x^2/4\nu}$ & relative error\\
\midrule
$400$ & $20$  & $1.00$ & $1.28$              & $0$ \\
$400$ & $50$  & $2.50$ & $4.77$              & $1.1\times10^{-13}$\\
$50$  & $20$  & $2.83$ & $7.39$              & $1.4\times10^{-14}$\\
$200$ & $50$  & $3.54$ & $2.28\times10^{1}$  & $8.5\times10^{-14}$\\
$400$ & $100$ & $5.00$ & $5.18\times10^{2}$  & $1.8\times10^{-12}$\\
$100$ & $50$  & $5.00$ & $5.18\times10^{2}$  & $2.2\times10^{-12}$\\
$200$ & $100$ & $7.07$ & $2.68\times10^{5}$  & $6.8\times10^{-7}$\\
\bottomrule
\end{tabular}
\caption{Branch A against \code{mpmath} at 50 digits. The error tracks $\eps M$. The last
row lies past the threshold $M\le10^3$ and is where the code hands over to branch B.}
\label{tab:seriesacc}
\end{table}

The loop terminates when $|t_j|\le\eps|S|$. Term counts are small: three terms at $x=0.01$,
twenty-four at $(\nu,x)=(400,50)$, forty-five at $(400,100)$.

It is worth being precise about the worst case, because it is not where one first looks.
Sweeping the whole admitted region $M\le10^3$, the largest relative error found is
$2.9\times10^{-11}$, at $(\nu,x)=(400,105)$ --- and $M$ there is not exceptional. What is
exceptional is $|S|\approx9.8\times10^{-4}$: the point sits near a zero of $J_\nu$, and the
error is $\eps M/|S|$, with $|S|$ in the \emph{denominator}. Absolute accuracy is untouched;
it is relative accuracy that degrades near a zero, as it must. The bound the code reports
there is $3.3\times10^{-11}$, so the diagnostic predicts its own worst case rather than being
surprised by it.

\subsection{The diagnostic that falls out}

$M$ is computed anyway inside the loop, so the bound $\eps M/|S|$ costs nothing to report.
The code returns it as \code{rel\_err}, and the basis reduces it over the evaluation points
to a per-column relative-accuracy bound. This is the cheap-bound-and-warn pattern the
geometry layer uses for \code{Domain.achievable\_prec}: compute a bound in microseconds,
warn only when the request exceeds it.

\section{Branch B: library evaluation, rescaled}

Where $M$ is too large the series is abandoned and \code{scipy.special.jv} is called
directly, then rescaled:
\[
  S=\sgn\big(J_\nu(x)\big)\exp\!\big(\ln|J_\nu(x)|-\sigma_J\big).
\]
The rescaling is done in logarithms rather than as $J\,e^{-\sigma}$ so that a large $\sigma$
cannot overflow on the way. A genuine zero of $J_\nu$ survives: $\ln 0=-\infty$ and
$e^{-\infty}=0$.

The derivative uses the symmetric recurrence
\[
  2J_\nu'(x)=J_{\nu-1}(x)-J_{\nu+1}(x),
\]
evaluated with two further \code{jv} calls. Two remarks on that choice. First, the
asymmetric form $J_\nu'=J_{\nu-1}-(\nu/x)J_\nu$ needs only one extra call but suffers
cancellation wherever the two terms are close, which the symmetric form does not. Second,
\code{scipy.special.jvp} is not used: it routes every call through a generic $n$-th
derivative formula with a Python loop and a binomial coefficient per term, and profiling a
realistic $600\times144$ basis block put it at $39\%$ of the whole per-$k$ evaluation ---
more than the Bessel functions it was calling.

Because those two extra calls triple the cost of the library branch, the code takes a flag:
when the caller wants values only --- which is every Dirichlet boundary, the common case ---
the derivative is not formed at all. Measured on that same block, $26.3$\,ms with the
derivative against $13.2$\,ms without.

\section{Branch C: the uniform (Debye) asymptotic}

Where the library underflows \emph{and} the series is untrustworthy, the code falls back on
the classical uniform asymptotic for large order (DLMF 10.19.6--10.19.7). With
\[
  z=\frac{x}{\nu},\qquad w=\sqrt{1-z^2},\qquad p=\frac1w,\qquad
  \eta = w+\ln\frac{z}{1+w} \quad (z<1),
\]
one has
\begin{equation}
  J_\nu(\nu z)\sim \frac{e^{\nu\eta}}{\sqrt{2\pi\nu}\,\sqrt{w}}\sum_{k\ge0}\frac{U_k(p)}{\nu^k},
  \qquad
  J_\nu'(\nu z)\sim \frac{\sqrt{w}\,e^{\nu\eta}}{\sqrt{2\pi\nu}\;z}\sum_{k\ge0}\frac{V_k(p)}{\nu^k},
  \label{eq:debye}
\end{equation}
with $U_0=V_0=1$ and
\[
\begin{aligned}
  U_1 &= \tfrac{1}{24}\left(3p-5p^3\right), &
  V_1 &= \tfrac{1}{24}\left(-9p+7p^3\right),\\
  U_2 &= \tfrac{1}{1152}\left(81p^2-462p^4+385p^6\right), &
  V_2 &= \tfrac{1}{1152}\left(-135p^2+594p^4-455p^6\right),\\
  U_3 &= \tfrac{1}{414720}\left(30375p^3-369603p^5+765765p^7-425425p^9\right), &&\\
  V_3 &= \tfrac{1}{414720}\left(-42525p^3+451737p^5-883575p^7+475475p^9\right). &&
\end{aligned}
\]
Reading off the amplitude in~\eqref{eq:debye}, the natural Debye scale is
\begin{equation}
  \sigma_D = \nu\eta - \tfrac12\ln\!\big(2\pi\nu w\big),
  \label{eq:sigmaD}
\end{equation}
and converting to the canonical scale of~\eqref{eq:factor} costs one exponential:
\[
  S = e^{\sigma_D-\sigma_J}\sum_k \frac{U_k(p)}{\nu^k},
  \qquad
  S' = \frac{w}{z}\,e^{\sigma_D-\sigma_J}\sum_k \frac{V_k(p)}{\nu^k}.
\]
Note that the polynomials are sparse in $p$ --- only every other power appears --- and
storing them as dense coefficient tuples invited an off-by-one that showed up as a
$2\times10^{-7}$ disagreement with the series. They are stored as explicit
(power, coefficient) pairs.

Four terms give at worst $7\times10^{-13}$ over the corner of parameter space where this
branch is actually reached ($\nu\gtrsim520$, $z\lesssim0.3$; measured against \code{mpmath}
at $(\nu,x)$ of $(520,100)$, $(600,130)$, $(600,180)$, $(700,200)$ and $(1000,300)$). It degrades only as $z\to1$, the Airy transition, which is
precisely where the library branch is healthy.

\section{Branch selection, and why the branches overlap}

The rule is:
\begin{center}
\begin{tabular}{ll}
\toprule
A & $M\le 10^3$, i.e.\ $x\lesssim 5.25\sqrt{\nu}$\\
B & otherwise\\
C & B returned exactly $0.0$\\
\bottomrule
\end{tabular}
\end{center}
The mask is taken over the whole $(\text{points}\times\text{columns})$ block, so each branch
is a single vectorised call.

Branch C selects \emph{itself}: rather than a threshold in $(\nu,x)$, the code runs branch B
and falls through wherever the library gave up. That is one fewer tuned constant, and it
cannot drift out of date if the library's own underflow edge moves.

The safety of the scheme rests on the branches overlapping rather than abutting. Sweeping
$\nu\in[0.2,600]$ against $x\in[0.5,100]$ on a $220\times220$ logarithmic grid and counting
cells where \code{jv} underflows \emph{and} $M>10^3$:
\[
  \#\{\text{cells in the gap}\} = 0 \quad\text{out of } 48\,400 .
\]
Loosening the series threshold to $M>10^2$ opens the first two gap cells, at $\nu\approx518$.
So the two principal branches are simultaneously valid over a wide strip, and branch C is a
safety net for $\nu$ beyond roughly $500$ rather than a workhorse. The overlap is also the
sharpest available test of the implementation: the two branches are compared against each
other throughout it, and that comparison fails first if either threshold is edited.

\section{The error floor is the scale, not the summation}
\label{sec:floor}

The bound $\eps M/|S|$ of \S3 is not the whole story, and the part it misses dominates in the
regime this note is about.

The scale $\sigma_J=\nu\ln(x/2)-\ln\Gamma(\nu+1)$ is assembled from two large numbers. At
$(\nu,x)=(400,50)$ they are $1287.6$ and $2000.5$. Each is computed to a relative accuracy of
about $\eps$, hence to an \emph{absolute} accuracy of about $\eps\times 2\times10^{3}\approx
2\times10^{-13}$ --- and an absolute error in a logarithm is a relative error in the value it
scales. The realistic bound is therefore
\begin{equation}
  \frac{\delta J}{J} \;\approx\; \eps\left(\frac{M}{|S|}
    + \big|\nu\ln(x/2)\big| + \big|\ln\Gamma(\nu+1)\big|\right).
  \label{eq:floor}
\end{equation}
Table~\ref{tab:floor} shows how badly the first term alone underestimates the truth.

\begin{table}[h]
\centering
\begin{tabular}{rrrrr}
\toprule
$\nu$ & $x$ & $\eps M/|S|$ & $\eps(|\nu\ln\frac{x}{2}|+|\ln\Gamma(\nu+1)|)$ & measured\\
\midrule
$400$ & $50$ & $1.1\times10^{-15}$ & $7.3\times10^{-13}$ & $1.1\times10^{-13}$\\
$400$ & $20$ & $2.9\times10^{-16}$ & $6.5\times10^{-13}$ & $0$\\
$100$ & $10$ & $2.9\times10^{-16}$ & $1.2\times10^{-13}$ & $2.8\times10^{-14}$\\
\bottomrule
\end{tabular}
\caption{The scale, not the summation, sets the floor. Omitting the middle column would
understate the error by two orders of magnitude, in the unsafe direction.}
\label{tab:floor}
\end{table}

\begin{remark}[The floor is mostly invisible to the solver]
Along a single basis column $\nu$ is fixed, so $\ln\Gamma(\nu+1)$ contributes a constant
multiplicative error to the whole column. The generalised singular value decomposition that
consumes these columns is invariant to column scaling, so it does not see that part at all;
what it sees is the summation term and the $r$-dependence of $\nu\ln(x/2)$. The code reports
the honest entrywise number from~\eqref{eq:floor} and says in its docstring that the bound is
conservative for a scale-invariant consumer.
\end{remark}

The practical consequence is a statement worth making once: \emph{the scaled kernel is a
$5\times10^{-13}$-relative object at $\nu\approx400$, not a $10^{-15}$ one.} Tests are written
against that number.

\section{The second kind}

Fundamental solutions use $Y_m(kr)$ at integer order, and the difficulty is the same one with
the sign reversed: $Y_m$ blows up as $x\to0$ where $J_\nu$ vanishes. The library overflows
rather than underflowing, at the last finite orders measured in
Table~\ref{tab:yoverflow}, and again the transition is a cliff at the format's ceiling rather
than a decay.

\begin{table}[h]
\centering
\begin{tabular}{lrrrrr}
\toprule
$x$ & $10^{-3}$ & $10^{-2}$ & $1$ & $5$ & $20$\\
\midrule
last finite $m$ & $65$ & $80$ & $149$ & $205$ & $294$\\
\bottomrule
\end{tabular}
\caption{The overflow edge of \code{scipy.special.yv}; the value there is about $10^{302}$.}
\label{tab:yoverflow}
\end{table}

The standard expansion for integer order $n\ge1$ is
\[
  Y_n(x) = -\frac{1}{\pi}\Big(\frac{2}{x}\Big)^{n}\sum_{j=0}^{n-1}\frac{(n-j-1)!}{j!}
           \Big(\frac{x^2}{4}\Big)^{j}
         + \frac{2}{\pi}\ln\frac{x}{2}\,J_n(x)
         - \frac{1}{\pi}\Big(\frac{x}{2}\Big)^{n}\!\!\sum_{j\ge0}\!\big[\psi(j{+}1){+}\psi(n{+}j{+}1)\big]
           \frac{(-x^2/4)^j}{j!\,(n+j)!}.
\]
Only the first sum carries the singular behaviour. Factoring out its $j=0$ term gives the
scale and the recurrence
\begin{equation}
  \sigma_Y = \ln\Gamma(n) - n\ln\frac{x}{2} - \ln\pi,
  \qquad a_0=1,\quad a_{j+1}=a_j\,\frac{x^2/4}{(n-j-1)(j+1)},
\end{equation}
with $S=-\sum_{j=0}^{n-1}a_j$ and, differentiating termwise exactly as in~\eqref{eq:deriv},
$x S' = -\sum_j (2j-n)\,a_j$. For $n=0$ the singularity is logarithmic rather than algebraic,
so no power-law scale applies and the code sets $\sigma_Y=0$ and evaluates directly; the
order-zero path also uses \code{y0} and the exact identity $Y_0'=-Y_1$, which is
substantially faster than the general-order routines.

The two dropped groups of terms are relatively small by the factor
\begin{equation}
  \frac{\text{dropped}}{\text{kept}} \;\sim\;
  \exp\!\Big(2n\ln\tfrac{x}{2}-\ln\Gamma(n+1)-\ln\Gamma(n)+\ln\pi\Big),
  \label{eq:ydrop}
\end{equation}
and the branch condition is exactly that this be below $10^{-16}$. As with $J_\nu$ the two
branches overlap: the truncated sum is poor only for $n\lesssim5$ with $x\gtrsim n$, which is
precisely where the library returns an $O(1)$ number with full accuracy.

\begin{remark}
Choosing the branch on $\sigma_Y$ rather than on~\eqref{eq:ydrop} looks equivalent and is not.
At $(n,x)=(5,0.5)$ the scale is already $e^{9}$, which sounds safe, while the dropped
logarithmic term is $9\times10^{-10}$ of the total. The ratio is the quantity that governs
the truncation, and it is the one to test.
\end{remark}

\section{The per-column scale}
\label{sec:column}

The pointwise scale $\sigma_J(\mu,kr)$ of \S2 makes each \emph{entry} representable. It does
not make a \emph{column} well scaled, because it varies with $r$ along the column. The basis
therefore divides each column $j$ by a second, purely analytic quantity $\sigma_j(k)$ that
depends on the column and on $k$ but on nothing about the evaluation points:
\[
  A[:,j] \;=\; \psi_j\,e^{-\sigma_j(k)},
  \qquad
  \psi_j(z)=J_\mu(k r)\sin(\mu\theta).
\]
In the evaluation this appears as the single factor $e^{\sigma_J(\mu,kr)-\sigma_j(k)}$, which
is where the two scales meet and why both must be canonical analytic numbers.

Point-set independence is what makes this compose with anything the solver does downstream.
If the generalised singular value decomposition further rescales columns to unit norm,
$\hat A = A\,\mathrm{diag}(1/d)$, then the coefficient of $\psi_j$ is
$\hat c_j\,e^{-\sigma_j}/d_j$: the two log-scales simply add, and neither layer needs to know
the other's number. (They are never exponentiated separately. At a $\pi/8$ corner with
$kR=10$ the $\sigma_j$ span about $1149$ nats, so $e^{\sigma_j}$ ranges over $10^{-499}$.)

\subsection{Construction}

Let $R$ be the largest radius at which the family will be evaluated --- the maximum of
$|z-\text{origin}|$ over $\Omega$, a fixed geometric constant.\footnote{Not the corner's
clearance, which bounds where the \emph{local expansion} is valid rather than where the
column is \emph{evaluated}. Using the clearance is an error of $(r_{\max}/R)^\mu$, which on
an L-shape with twelve degrees per corner was measured at $5.7\times10^{10}$.} Two regimes
compete. For $kR<\mu$ the maximum of $|J_\mu|$ over $r\le R$ sits at $r=R$ and is given by
the Debye amplitude~\eqref{eq:sigmaD}; past $kR=\mu$ the function is oscillatory and the
relevant bound is the global maximum over all arguments,
\[
  \max_{x\ge0} J_\mu(x) \;\approx\; C_A\,(\mu+\delta)^{-1/3},
  \qquad C_A=0.6748,\;\delta=0.3,
\]
the $C_A$ being the Airy constant and the shift $\delta$ fitted so that the same expression
is also right at small order (measured ratio to the true maximum in $[0.934,1.000]$ for
$0\le\mu\le352$). Writing $a$ for the log of the first and $b$ for the log of the second, the
scale is a smooth minimum,
\begin{equation}
  \sigma_j \;=\; -\frac{1}{p}\ln\!\big(e^{-pa}+e^{-pb}\big),\qquad p=4 .
  \label{eq:softmin}
\end{equation}

\subsection{Three corrections, each forced by a measurement}

\paragraph{The blend must be smooth.} A hard $\min$ has a discontinuous derivative at the
crossover, and $\sigma_j(k)$ sits inside the objective the eigensolver fits a parabola to. The
soft minimum costs a factor $2^{1/p}=1.19$ in the column maximum where the two branches meet,
which is the price of that.

\paragraph{The amplitude must be floored.} The Debye amplitude
$-\tfrac12\ln(2\pi\mu w)$ diverges as $w\to0$, which is exactly the crossover. Measured as a
second difference of $\sigma_j$ in $k$ across $kR=\mu$ at $\mu=20$:
\begin{center}
\begin{tabular}{lr}
\toprule
unfloored amplitude, $p=1$ & $8.4\times10^{5}$\\
unfloored amplitude, hard $\min$ & $1.9\times10^{3}$\\
floored amplitude, $p=4$ (as shipped) & $6.5$\\
\bottomrule
\end{tabular}
\end{center}
The middle row is the instructive one: \emph{unfloored, the smooth blend is worse than
$\min$}, because $\min$ clips the divergence while $\ln(e^{-a}+e^{-b})$ follows it in. The
floor is the load-bearing correction; the soft blend removes the kink that remains.

The floor is placed at the Airy transition width, where the asymptotic ceases to be valid
anyway, and its constant is \emph{solved} rather than tuned: requiring the floored amplitude
to equal the global maximum exactly at $z=1$ (where $\eta=0$) gives
\[
  -\tfrac12\ln\big(2\pi\mu\,w_{\text{floor}}\big) = \ln\!\big(C_A(\mu+\delta)^{-1/3}\big)
  \quad\Longrightarrow\quad
  w_{\text{floor}} = \frac{C_A^{-2}}{2\pi}\cdot\frac{(\mu+\delta)^{2/3}}{\mu},
\]
with $C_A^{-2}/(2\pi)=0.349518$. So the envelope and the cap \emph{meet} at the crossover
rather than merely crossing.

\paragraph{The positive part must be smooth too.} Writing $q=1-z^2$, the floored amplitude
uses $\big(q_+^2+w_{\text{floor}}^4\big)^{1/4}$. Past the crossover $q$ turns negative, and
using $q^2$ unmodified lets the amplitude grow again, driving the column maximum to $27.7$.
Clamping with $\max(q,0)$ fixes the magnitude but leaves a jump in the second derivative;
the code uses the $C^\infty$ form
\[
  q_+ = \tfrac12\Big(q+\sqrt{q^2+w_{\text{floor}}^4}\Big),
\]
which has the same two limits.

\subsection{What it achieves}

Over $\mu$ from $0.05$ to $320$ and $kR$ from $0.5$ to $300$, the maximum of each scaled
column lies in $[0.802,\,1.185]$ --- a spread of $1.48$ --- against raw columns spanning
hundreds of orders of magnitude on the same grid. The second difference of $\sigma_j$ in $k$
is at most $20.8$ at $\mu=2$, $3.1$ at $\mu=20$ and $0.5$ at $\mu=200$.

Finally, $\sigma_j$ depends on $k$ and $R$ only through the product $kR$. Under the geometric
normalisation $\Omega\mapsto s\Omega$, for which $\lambda\mapsto s^{-2}\lambda$ and hence
$k\mapsto k/s$, the product is unchanged and every column scale is invariant \emph{exactly},
bit for bit. That is asserted as an equality rather than to a tolerance: if it ever needs a
tolerance, the formula has grown a term that should not be there.

\section{Summary}

\begin{center}
\begin{tabular}{lll}
\toprule
& what it computes & accuracy\\
\midrule
A, series & $S$ and $S'$ from one loop & $\eps M/|S|$; $\le 3\times10^{-11}$ over $M\le10^{3}$\\
B, library & $\code{jv}$, rescaled in logs & $\le 2\times10^{-13}$ up to $\nu\approx400$\\
C, Debye & four-term uniform asymptotic & $\le 7\times10^{-13}$ where reachable\\
\midrule
floor & $\eps\,|\sigma_J|$, from \S\ref{sec:floor} & $\sim10^{-13}$ at $\nu=400$\\
\bottomrule
\end{tabular}
\end{center}

\noindent
Away from the zeros of $J_\nu$ --- which is almost everywhere --- the series delivers
$10^{-13}$. The $3\times10^{-11}$ is the price of asking for \emph{relative} accuracy at a
point where the function is a thousand times smaller than its own series terms, and the
reported bound says so in advance.

The three branches overlap with no gap over the parameter range the solver uses; the
factorisation makes $|S|\le1$ a checkable invariant; the derivative shares the value's loop
and its scale; and a second, purely analytic per-column scale turns a $500$-decade range into
a $1.5$-fold one without ever consulting the evaluation points.

\section*{References}

\begin{itemize}
\item NIST Digital Library of Mathematical Functions, \S10.2 (series), \S10.6 (recurrences),
      \S10.19 (uniform asymptotics for large order). \url{https://dlmf.nist.gov/}
\item G.\,N. Watson, \emph{A Treatise on the Theory of Bessel Functions}, 2nd ed., \S3.31
      for the bound of Proposition 1.
\item F.\,W.\,J. Olver, \emph{Asymptotics and Special Functions}, ch.~12, for the $U_k$ and
      $V_k$ polynomials and the range of validity of~\eqref{eq:debye}.
\end{itemize}

\end{document}
